\documentclass[11pt]{article}
\usepackage{amsmath,amssymb}
\usepackage[margin=1in]{geometry}
\usepackage{hyperref}
\emergencystretch=3em

\title{The explicit coefficient series and the phase-derivative dictionary (Taylor-tree Stage 5)}
\author{Claude, with Daniel Murfet}
\date{2026-09-07}

\begin{document}
\maketitle
\sloppy

\begin{abstract}
Units 250--254 close review v21's qualification (ii): the family coefficients $A_{\mu,j}(\xi,\eta)$ of
Stage~4, defined as limits of polynomial truncation coefficients, are identified with the paper's
explicit absolutely convergent series. With $J$ the constant-free part of the coefficient family of
$\xi$ and $c*e$ the Cauchy product $(c*e)_\gamma=\sum_{\alpha\le\gamma}c_\alpha e_{\gamma-\alpha}$,
\begin{align*}
A_{\mu,j}(\xi,\eta)&=\sum_{p\ge0}\frac{\beta^p}{p!}\,T_{\mu,j,p}\bigl(c^\eta*J^{*p}\bigr),\\
T_{\mu,j,p}(f)&=\prod_i\frac1{2k_i}\sum_\gamma f_\gamma\sum_{q=j}^{d-1}c_{\mu,q}(u^{h+\gamma})\binom qj
\int_0^\infty t^{\mu-1}(-\log t)^{q-j}(\sqrt t)^pe^{-\beta t+\beta\sqrt ta}\,dt
\end{align*}
(Headline~XXIX, \texttt{familySpectralCoeff\_eq\_series}, $\mu>0$), where $(c^\eta*J^{*p})_\gamma=
\sum_{\mathbf m+\mathbf n=\gamma}\frac{\eta_{\mathbf m}}{\mathbf m!}\xi_{\mathbf n,p}$ is the paper's
collected coefficient (eq:flucttreeterms) once the families are normalised Taylor coefficients. The
phase half of the derivative dictionary is a theorem: $\partial_a^pS_\mu(a)=\beta^p\int_0^\infty
t^{\mu-1}(\sqrt t)^pe^{-\beta t+\beta\sqrt ta}\,dt$ (unit 250). Review v22: qualified pass, no
mathematical defect. No \texttt{sorry} and no additional \texttt{axiom} declarations; 9130 jobs.
\end{abstract}

\section{Architecture}
\paragraph{Cauchy product without an antidiagonal instance (u251).} Mathlib has no
\texttt{HasAntidiagonal} on $\mathrm{Fin}\,d\to\mathbb N$; the convolution is a finite sum over
$\mathrm{Iic}\,\gamma$ and the double series over pairs is regrouped along $\gamma=\alpha+\delta$ by the
equivalence $(\alpha,\delta)\mapsto\langle\alpha+\delta,\alpha\rangle$ and \texttt{HasSum.sigma}, giving
$\mathrm{eval}(c*e)=\mathrm{eval}\,c\cdot\mathrm{eval}\,e$ on the cube and $\mathrm{mass}(c*e)\le
\mathrm{mass}\,c\cdot\mathrm{mass}\,e$.
\paragraph{Collected coefficients (u252).} A monomial list collects to a finitely supported family;
list product collects to convolution, powers to convolution powers, the fluctuation part to the
constant-free family and box truncations to restrictions. Every Stage-3 list sum
$\sum_{(\gamma,c)\in P}c\,f(\gamma)$ equals $\sum_\gamma(\mathrm{coeffFn}\,P)_\gamma f(\gamma)$.
\paragraph{Kernel functional (u253).} $\mathrm{coeffTerm}_p(P)=T_p(\mathrm{coeffFn}\,P)$ with
$|T_p f|\le K_kD\,M_{\mu,d-1,p}(a)\,\mathrm{mass}f$: the coefficient functional is $\ell^1$-continuous,
uniformly in the monomial by the Stage-3 budget.
\paragraph{Identification (u254).} The truncation coefficients are the same series for the truncated
families; each term converges as the box grows (ℓ$^1$ continuity, the family power-difference estimate
$\mathrm{mass}(c^{*p}-c'^{*p})\le pB^{p-1}\mathrm{mass}(c-c')$, tail masses $\to0$), and the series
passes to the limit by dominated convergence with the Tonelli majorant
$\sum_p\frac{(\beta B)^p}{p!}M_{\mu,n,p}(a)=M_{\mu,n}(a+B)$. Uniqueness of limits identifies the
Stage-4 coefficient with the explicit series.
\paragraph{Phase derivatives (u250).} Differentiation under the integral
(\texttt{hasDerivAt\_integral\_of\_dominated\_loc\_of\_deriv\_le}) with the log majorant at phase
parameter $|a|+1$ gives $\partial_a\,\mathrm{fluctMoment}=\beta\,\mathrm{fluctMoment}$ at the next phase
order, hence $\partial_a^pS_\mu=\beta^pS_{\mu+p/2}$ (the paper's displayed formula omits $\beta^p$).

\section{Status of the programme}
Proved: the Taylor-tree expansion and remainder on general boxes for phase and amplitude given by
absolutely summable power series (Headlines~XXVII--XXVIII), with the paper's exponent set, logarithmic
indexing and cutoff-independent coefficients equal to the paper's explicit series (XXIX). Not proved:
the Cauchy-estimate bridge from holomorphy on a polydisc to summability of the Taylor coefficients (no
several-variable Cauchy estimates in Mathlib; Astra \#29 estimates 15--25 units and does not budget it),
and the spectral half $(-\partial_\mu)^i$ of the derivative dictionary (next unit).

\section{Lessons}
Dot notation on hypotheses of a \texttt{def}-alias type (\texttt{AbsSummable}) resolves to lemmas in
that namespace first; call \texttt{Summable.add} explicitly. \texttt{conv} already existed in the
project namespace: qualify as \texttt{CoeffFamily.conv}. \texttt{Pi.sub\_apply} before
\texttt{Finset.sum\_sub\_distrib}; \texttt{List.filter\_cons\_of\_pos/neg} for filters;
\texttt{Finset.sum\_coe\_sort} for fibre sums over a Finset coercion; the pi-type order instances
live in \texttt{Mathlib.Data.Pi.Interval} and \texttt{Mathlib.Algebra.Order.Sub.Prod}.
\end{document}
